% =====================================================================
% Related Work section for the SE2RoadNet paper (ICSE 2027 target).
% Self-contained: \section + a \thebibliography block using readable
% cite keys. If merging into the RoadFury paper.tex (which already has
% b1,b2,b6,b7,b9,b10,b11), map: sensodat->b7, sdcscissor->b1, saner->b2,
% comp26->b8, focal->b9, gcn->b10, transformer->b11, swa->b6, and drop
% the duplicate \bibitem lines.
% Compile with XeLaTeX x2 (see memory: latex-build-xelatex).
% =====================================================================

\section{Related Work}
\label{sec:related}

Prioritizing simulation-based tests for self-driving cars (SDCs) has been
attacked from three angles: (i)~\emph{search-based / diversity} methods that
reorder tests to maximize behavioral spread, (ii)~\emph{feature-based machine
learning} that classifies each road from hand-engineered statistics, and
(iii)~\emph{deep learning} that scores the road geometry end-to-end. All three
raise the average percentage of faults detected (APFD~\cite{apfd,yooharman})
over random ordering, yet none offers a \emph{guarantee} that the score is
stable under the two transformations that leave a road's physics unchanged:
rigid rotation/translation of the road in the plane (the group $SE(2)$) and a
change of point-sampling resolution. We review each line, then the invariance
theory we build on, and position SE2RoadNet as the first SDC prioritizer whose
invariance is architectural rather than augmentation-based.

\subsection{Search-based and diversity-driven prioritization}

SDC-Prioritizer~\cite{sdcprio} treats prioritization as black-box search: a
chromosome $\tau=\langle t_1,\dots,t_n\rangle$ is an execution order, and a
genetic algorithm evolves it so that the most diverse, cheapest tests run
first. The single-objective variant maximizes a cost-discounted diversity
score,
\begin{equation}
f(\tau)=\sum_{i=2}^{n}\frac{\mathit{distance}(t_i,t_{i-1})}{\mathit{cost}(t_i)\times i},
\label{eq:sosdc}
\end{equation}
where $\mathit{distance}$ is the Euclidean distance between two tests'
$z$-score-normalized static road-feature vectors and the position multiplier
$i$ pushes high-value tests early. The multi-objective variant
(MO-SDC-Prioritizer) splits this into a diversity objective
$\max f_1=\sum_{i\ge 2}\mathit{distance}(t_i,t_{i-1})/i$ and a cost objective
$\min f_2=\sum_{i\ge 1}\mathit{cost}(t_i)/i$, and returns an NSGA-II Pareto
front. The approach is fully static and cheap ($<0.5\%$ overhead) and
significantly beats random and greedy baselines. Its limitation is
foundational for our story: diversity is only a \emph{surrogate} for
fault-revealing power---the method never learns the map from road shape to
failure probability---and its Euclidean distance over global features
(direct distance, absolute turn angles, pivot radii) is neither
rotation-invariant in any principled sense nor invariant to the polyline
sampling rate.

\subsection{Feature-based machine learning}

SDC-Scissor~\cite{sdcscissor,saner} predicts, before any simulation, whether a
road is ``unsafe'' by feeding hand-engineered static descriptors---turn counts,
per-segment turn-angle and pivot-radius statistics, road length, and a
Shapely-polygon ``road diversity'' area---into off-the-shelf classifiers
(Logistic Regression, Random Forest, Gradient Boosting). It reports
cost-effectiveness $\mathrm{CE}=\#\text{failing}/\#\text{passing}$ up to $4.0$
and avoids $\sim\!50\%$ of unnecessary executions. Two structural weaknesses
motivate our design. First, its counting features (number of left/right turns,
number of straight segments) are defined on a chosen \emph{segmentation} of the
spline, so they change when the road is resampled---there is no resolution
guarantee. Second, it collapses the road into a fixed bag of \emph{global}
scalar summaries, discarding the ordered curvature-versus-arclength profile
(\emph{where} along the road a sharp curvature occurs), and performs binary
selection rather than fine-grained ranking.

\subsection{Deep sequence and image models}

The strongest prior results come from deep models that consume the road
geometry directly. ITS4SDC~\cite{its4sdc} represents a road as a sequence of
per-segment $(\theta_k,\ell_k)$ pairs---relative turning angle and segment
length---and classifies it with a single bidirectional LSTM (220 cells),
reaching $F_1=0.89$ on the ICST~2025 competition set. It sets the first
segment's angle to $0$ and stores successive \emph{angle differences}
$\theta_k=\alpha_k-\alpha_{k-1}$; this cancels the arbitrary global heading but
is a hand-crafted normalization, \emph{not} an exact invariance, and the LSTM's
fixed input length ($N{=}197$) makes it sampling-rate dependent (the paper
itself lists fixed input dimension as a drawback). Our own prior baseline,
RoadFury~\cite{roadfury}, is a Pre-LN Transformer~\cite{transformer,preln}
over a $10$-channel road representation stabilized with SWA~\cite{swa}; it wins
the ICST~2026 competition at APFD~$\approx 0.807$ but has two blind spots we
close here---three of its channels ($\sin\theta,\cos\theta,\theta$) are
frame-dependent, and it uses absolute sinusoidal positional
encoding~\cite{transformer}, so both a rotated road and a resampled road move
its score. Graph- and image-based deep models have also been explored for road
scenarios, but they share the same limitation: none constrains the score to be
invariant under road rotation or resampling---the property we make exact.

\subsection{Benchmark and metric}

We evaluate on SensoDat~\cite{sensodat}, a public set of $32{,}580$ executed
BeamNG.tech SDC tests produced by three SBST generators (Frenetic, FreneticV,
AmbieGen) with PASS/FAIL labels from an out-of-bound oracle. Ranking quality is
measured by APFD~\cite{apfd,yooharman},
\begin{equation}
\mathrm{APFD}=1-\frac{\sum_{i=1}^{m}TF_i}{n\,m}+\frac{1}{2n},
\label{eq:apfd}
\end{equation}
where $n$ is the number of tests, $m$ the number of faults, and $TF_i$ the rank
of the first test exposing fault $i$; higher is better and random ordering
yields $\approx 0.5$.

\subsection{Invariance, and the components we build on}

Our guarantee is grounded in the theory of group-equivariant
networks~\cite{gcnn}. A map $\Phi$ is \emph{equivariant} to a group $G$ if
\begin{equation}
\Phi(T_g\,x)=T'_g\,\Phi(x)\qquad\forall g\in G,
\label{eq:equivariance}
\end{equation}
and \emph{invariance} is the special case $T'_g=\mathrm{Id}$, i.e.
$\Phi(T_g x)=\Phi(x)$. Cohen and Welling realize equivariance by convolving
over a \emph{discrete} group ($p4$ rotations, $p4m$ rotation-reflections), which
gives exact guarantees only for $90^{\circ}$ steps on a fixed pixel grid---no
continuous $SE(2)$ invariance and no resolution invariance. SE2RoadNet instead
attains \emph{exact, continuous} $SE(2)$ invariance without group averaging, by
(a)~encoding each road with intrinsic, coordinate-free features
(signed curvature $\kappa$ and its arclength derivatives), which are invariant
under any rigid motion by the fundamental theorem of planar
curves~\cite{frenet}, and (b)~injecting position through a relative-arclength
attention bias built with Random Fourier Features~\cite{rff}. The latter is the
scalar case of the RFF map $z(\Delta)=\sqrt{2/D}\,[\cos(\omega_1\Delta+b_1),\dots]$
whose inner product approximates a shift-invariant kernel $k(\Delta)$
(Bochner's theorem); because it depends only on the arclength \emph{difference}
$\Delta s=s_i-s_j$, the bias is invariant to the arclength origin and to
reparametrization. The remaining ingredients are standard: multi-head
self-attention~\cite{transformer} as the backbone, focal loss~\cite{focal},
\begin{equation}
\mathrm{FL}(p_t)=-(1-p_t)^{\gamma}\log(p_t),
\label{eq:focal}
\end{equation}
to counter the imbalanced FAIL class ($p_t$ is the predicted probability of the
true class and $\gamma$ down-weights easy examples), and SWA~\cite{swa} for
flat-minima generalization.

\subsection{Positioning}

Table~\ref{tab:related} summarizes the landscape along two axes: predictive
APFD, on which the field has converged near $0.80$, and the invariance
guarantee, on which \emph{every} prior method is empty. SE2RoadNet's
contribution is not another fraction of APFD but a new axis---exact $SE(2)$
rotation invariance (measured $\Delta\mathrm{APFD}=0.0000$ across six rotation
angles) with resolution invariance and curvature-monotone auditability---at
APFD parity with the best baseline and a clear AUC gain. In a
safety-critical setting, a provable guarantee is worth more than a few
thousandths of APFD.

% ---------------------------------------------------------------------
% Optional comparison table (fill APFD column from your harness; the
% "Invariant" column is the headline contrast).
% ---------------------------------------------------------------------
\begin{table}[t]
\caption{Prior SDC prioritization methods vs.\ SE2RoadNet. The rightmost
column---an architectural $SE(2)$/resolution guarantee---is the axis no prior
method occupies.}
\label{tab:related}
\centering{\small
\setlength{\tabcolsep}{4pt}
\begin{tabular}{llcc}
\toprule
\textbf{Method} & \textbf{Family} & \textbf{APFD} & \textbf{Invariant} \\
\midrule
Random                       & baseline        & $\approx 0.50$ & \ding{55} \\
SO-SDC-Prioritizer~\cite{sdcprio} & search     & $\approx 0.77$ & \ding{55} \\
Greedy diversity~\cite{sdcprio}   & search     & $\approx 0.80$ & \ding{55} \\
ITS4SDC~\cite{its4sdc}       & deep (bi-LSTM)  & $\approx 0.78$ & \ding{55} \\
RoadFury~\cite{roadfury}     & deep (Transf.)  & $\approx 0.81$ & \ding{55} \\
\textbf{SE2RoadNet (ours)}   & deep (Transf.)  & $\approx 0.81$ & \textbf{\ding{51}} \\
\bottomrule
\end{tabular}}
\end{table}

% =====================================================================
% Bibliography for this section (readable keys). Merge into the paper's
% main \thebibliography, deduplicating shared entries.
% Requires \usepackage{pifont} for \ding{51}/\ding{55} in Table above.
% =====================================================================
\begin{thebibliography}{99}\small

\bibitem{sensodat} C.~Birchler, C.~Rohrbach, T.~Kehrer, and S.~Panichella,
``SensoDat: Simulation-based sensor dataset of self-driving cars,'' in
\textit{Proc. IEEE/ACM MSR}, 2024. arXiv:2401.09808.

\bibitem{sdcprio} C.~Birchler, S.~Khatiri, P.~Derakhshanfar, S.~Panichella, and
A.~Panichella, ``Single and multi-objective test cases prioritization for
self-driving cars in virtual environments,'' \textit{ACM Trans. Softw. Eng.
Methodol. (TOSEM)}, vol.~32, no.~2, 2023. arXiv:2107.09614.

\bibitem{sdcscissor} C.~Birchler, S.~Khatiri, B.~Bosshard, A.~Gambi, and
S.~Panichella, ``Machine learning-based test selection for simulation-based
testing of self-driving cars software,'' \textit{Empir. Softw. Eng.}, vol.~28,
no.~71, 2023. arXiv:2212.04769.

\bibitem{saner} C.~Birchler, N.~Ganz, S.~Khatiri, A.~Gambi, and S.~Panichella,
``Cost-effective simulation-based test selection in self-driving cars software
with SDC-Scissor,'' in \textit{Proc. IEEE SANER}, 2022, pp.~386--390.

\bibitem{its4sdc} A.~G\"{u}ll\"{u}, F.~A.~Shah, and D.~Pfahl, ``An LSTM-based
test selection method for self-driving cars,'' in \textit{Proc. IEEE ICST Tool
Competition (SDC Testing Track)}, 2025. arXiv:2501.03881.

\bibitem{roadfury} [Authors], ``RoadFury: A Transformer with stochastic weight
averaging for SDC test prioritization,'' ICST~2026 SDC Testing Competition,
2026.

\bibitem{transformer} A.~Vaswani \textit{et al.}, ``Attention is all you
need,'' in \textit{Proc. NeurIPS}, 2017.

\bibitem{preln} R.~Xiong \textit{et al.}, ``On layer normalization in the
Transformer architecture,'' in \textit{Proc. ICML}, 2020.

\bibitem{focal} T.-Y.~Lin, P.~Goyal, R.~Girshick, K.~He, and P.~Doll\'{a}r,
``Focal loss for dense object detection,'' in \textit{Proc. IEEE ICCV}, 2017,
pp.~2980--2988.

\bibitem{swa} P.~Izmailov, D.~Podoprikhin, T.~Garipov, D.~Vetrov, and
A.~G.~Wilson, ``Averaging weights leads to wider optima and better
generalization,'' in \textit{Proc. UAI}, 2018.

\bibitem{gcnn} T.~S.~Cohen and M.~Welling, ``Group equivariant convolutional
networks,'' in \textit{Proc. ICML}, 2016, pp.~2990--2999.

\bibitem{rff} A.~Rahimi and B.~Recht, ``Random features for large-scale kernel
machines,'' in \textit{Proc. NeurIPS (NIPS)}, 2007.

\bibitem{apfd} S.~Elbaum, A.~G.~Malishevsky, and G.~Rothermel, ``Test case
prioritization: A family of empirical studies,'' \textit{IEEE Trans. Softw.
Eng.}, vol.~28, no.~2, pp.~159--182, 2002.

\bibitem{yooharman} S.~Yoo and M.~Harman, ``Regression testing minimization,
selection and prioritization: A survey,'' \textit{Softw. Test. Verif. Reliab.},
vol.~22, no.~2, pp.~67--120, 2012.

\bibitem{frenet} M.~P.~do~Carmo, \textit{Differential Geometry of Curves and
Surfaces}. Prentice-Hall, 1976.

\bibitem{comp26} P.~Aryan, T.~Fulcini, L.~Starace, C.~Birchler, and
S.~Panichella, ``ICST tool competition 2026---SDC testing track,'' in
\textit{Proc. IEEE ICST}, 2026.

\end{thebibliography}
